\chapter{Тестирование и оценка системы}

\section{Стратегия тестирования}

Стратегия тестирования системы Guide Helper предусматривает несколько уровней проверки качества программного обеспечения в соответствии с методологией, описанной в~\cite{kulikov2020testing,sommerville2016,pressman2014}.

Модульное тестирование (unit testing) направлено на проверку корректности отдельных функций и методов в изоляции от внешних зависимостей. Тесты написаны на языке Rust с использованием встроенной системы тестирования (\texttt{\#[cfg(test)]}). В модульных тестах внешние зависимости (база данных, NATS, внешние API) не используются: тесты оперируют только доменными объектами и чистыми функциями.

Интеграционное тестирование проверяет взаимодействие компонентов системы: HTTP-эндпоинтов, промежуточного слоя аутентификации, маршрутизатора Axum. Для интеграционных тестов применяется тестовый экземпляр приложения с настоящей PostgreSQL-базой данных в контейнере Docker.

Функциональное тестирование проверяет выполнение сформулированных требований к системе на уровне пользовательских сценариев. Тестирование функциональных требований выполнялось вручную по сценариям критического пути для трёх ключевых ролей пользователей.

Нагрузочное тестирование проводилось для проверки соответствия нефункциональным требованиям по производительности: время отклика API на 95-м перцентиле при 100 одновременных запросах не должно превышать 500 миллисекунд.

\section{Модульное тестирование}

Для сервиса маршрутов разработано 115 модульных тестов, покрывающих следующие модули:
\begin{itemize}
  \item доменная модель (\texttt{Route}, \texttt{RoutePoint}, \texttt{PhotoData})~--- 12 тестов;
  \item JWT-сервис~--- 8 тестов, включая проверку истекших и невалидных токенов;
  \item задача обработки фото (\texttt{PhotoProcessTask})~--- 6 тестов;
  \item бизнес-логика маршрутов, комментариев, лайков, рейтингов~--- 34 теста;
  \item импорт GeoJSON~--- 11 тестов на различные форматы и граничные случаи;
  \item категории и уведомления~--- 14 тестов;
  \item HTTP-обработчики (с тестовым приложением Axum)~--- 18 тестов;
  \item AI-чат и rate limiting~--- 12 тестов.
\end{itemize}

В листинге~\ref{lst:unit-test-example} приведён пример модульного теста, проверяющего обратную совместимость десериализации поля \texttt{photo}: система должна принимать как строку в формате data URL (старый формат), так и структуру \texttt{PhotoData} (новый формат).

\begin{lstlisting}[language=Rust, caption={Тест обратной совместимости десериализации фото}, label={lst:unit-test-example}]
#[test]
fn test_backward_compat_plain_string_photo() {
    // old format: photo field is a plain string
    let json = r#"{"lat":55.0,"lng":37.0,"name":null,
"segment_mode":null,"photo":"data:image/png;base64,abc"}"#;

    let point: RoutePoint = serde_json::from_str(json).unwrap();

    let photo = point.photo.unwrap();
    assert_eq!(photo.original, "data:image/png;base64,abc");
    assert_eq!(photo.status, PhotoStatus::Pending);
    assert!(photo.thumbnail_url.is_none());
}

#[test]
fn test_photo_data_struct_deserialization() {
    // new format: photo field is a PhotoData object
    let json = r#"{"lat":55.0,"lng":37.0,"name":null,
"segment_mode":null,"photo":{"original":"data:image/png;base64,abc",
"thumbnail_url":"/photos/thumb.jpg","status":"done"}}"#;

    let point: RoutePoint = serde_json::from_str(json).unwrap();
    let photo = point.photo.unwrap();
    assert_eq!(photo.status, PhotoStatus::Done);
    assert_eq!(photo.thumbnail_url,
               Some("/photos/thumb.jpg".to_string()));
}
\end{lstlisting}

Все тесты выполняются автоматически в CI/CD-пайплайне при каждом коммите. Команда \texttt{cargo test} запускает весь набор тестов; при падении хотя бы одного теста сборка Docker-образов не выполняется.

\section{Функциональное тестирование}

Функциональное тестирование системы проводилось по методу критического пути пользователей~\cite{kulikov2020testing}. Для каждой из трёх ключевых ролей системы определён критический путь~--- минимальная последовательность действий, необходимая для достижения основной цели роли. Требования к системе сформулированы в соответствии с методикой С.~Куликова (раздел~2.2): каждое требование содержит идентификатор, наименование, описание и критерий тестирования. В данном разделе роли <<гид>> и <<турист>> рассматриваются как сценарии использования системы в рамках базовой системной роли \texttt{user}, а роль <<администратор>> соответствует системной роли \texttt{admin}.

\subsection{Сценарий для роли «Гид/инструктор»}

Основная цель роли: создать маршрут с геопривязанными фотографиями и опубликовать ссылку для клиентов. Результаты тестирования критического пути представлены в таблице~\ref{tab:test-guide}.

\begin{table}[H]
\caption{Результаты тестирования критического пути роли «Гид»}
\label{tab:test-guide}
\begin{tabularx}{\textwidth}{|c|X|X|c|c|}
\hline
\textbf{№} & \textbf{Действие} & \textbf{Ожидаемый результат} & \textbf{Тр-е} & \textbf{Рез.} \\
\hline
1 & Регистрация по email и паролю & Учётная запись создана, пользователь авторизован & FR-01 & \checkmark \\
\hline
2 & Создать маршрут, добавить 3 точки на карте & Маршрут сохранён, точки отображены & FR-06 & \checkmark \\
\hline
3 & Прикрепить JPEG с GPS к точке №2 & Статус точки «загрузка», фото в очереди & FR-13 & \checkmark \\
\hline
4 & Дождаться уведомления WebSocket & Миниатюра появилась на маркере карты & FR-14, FR-15 & \checkmark \\
\hline
5 & Проверить статистику маршрута & Расстояние, набор высоты, сложность рассчитаны & FR-12 & \checkmark \\
\hline
6 & Включить публичный доступ & Получена ссылка \texttt{/shared/\{token\}} & FR-09 & \checkmark \\
\hline
7 & Открыть ссылку без авторизации & Карта с маршрутом и фото доступна & FR-09 & \checkmark \\
\hline
\end{tabularx}
\end{table}

\subsection{Сценарий для роли «Турист»}

Основная цель роли: найти интересный маршрут и оставить социальную активность (оценку и комментарий). Результаты тестирования представлены в таблице~\ref{tab:test-tourist}.

\begin{table}[H]
\caption{Результаты тестирования критического пути роли «Турист»}
\label{tab:test-tourist}
\begin{tabularx}{\textwidth}{|c|X|X|c|c|}
\hline
\textbf{№} & \textbf{Действие} & \textbf{Ожидаемый результат} & \textbf{Тр-е} & \textbf{Рез.} \\
\hline
1 & Открыть каталог без входа & Список публичных маршрутов с пагинацией & FR-22 & \checkmark \\
\hline
2 & Ввести текст в поиск & Маршруты отфильтрованы по названию & FR-23 & \checkmark \\
\hline
3 & Открыть маршрут, кликнуть на маркер с фото & Всплывающее окно с фотографией & FR-09 & \checkmark \\
\hline
4 & Поставить оценку 5/5 & Средний рейтинг обновлён & FR-18 & \checkmark \\
\hline
5 & Написать комментарий & Комментарий в списке, автор получил уведомление & FR-16, FR-19 & \checkmark \\
\hline
6 & Добавить в закладки & Маршрут появился в разделе «Закладки» & FR-20 & \checkmark \\
\hline
\end{tabularx}
\end{table}

\subsection{Сценарий для роли «Администратор»}

Основная цель роли: управление контентом и пользователями платформы. Результаты тестирования представлены в таблице~\ref{tab:test-admin}.

\begin{table}[H]
\caption{Результаты тестирования критического пути роли «Администратор»}
\label{tab:test-admin}
\begin{tabularx}{\textwidth}{|c|X|X|c|c|}
\hline
\textbf{№} & \textbf{Действие} & \textbf{Ожидаемый результат} & \textbf{Тр-е} & \textbf{Рез.} \\
\hline
1 & Войти как admin & Навигация содержит «Администрирование» & FR-05 & \checkmark \\
\hline
2 & Список пользователей с ролями & Таблица всех пользователей & FR-24 & \checkmark \\
\hline
3 & Сменить роль пользователя & Роль изменена, права обновлены & FR-24 & \checkmark \\
\hline
4 & Удалить маршрут-нарушитель & Маршрут исчез из каталога & FR-26 & \checkmark \\
\hline
5 & Создать категорию «Горные» & Категория в фильтрах каталога & FR-25 & \checkmark \\
\hline
6 & Открыть статистику & Счётчики пользователей, маршрутов, фото & FR-27 & \checkmark \\
\hline
\end{tabularx}
\end{table}

\subsection{Статус реализации требований}

По результатам функционального тестирования, из 37 функциональных требований реализовано~33 (89\%). Техническое задание на систему оформлено в соответствии с~\cite{gost19201}, программа испытаний~--- в соответствии с~\cite{gost19301}. В таблице~\ref{tab:req-status} перечислены требования, находящиеся в разработке.

\begin{table}[H]
\caption{Требования, находящиеся в разработке}
\label{tab:req-status}
\begin{tabularx}{\textwidth}{|l|X|l|}
\hline
\textbf{ID} & \textbf{Требование} & \textbf{Статус} \\
\hline
FR-31 & AI-генерация описания по фотографиям & В разработке \\
\hline
FR-33 & QR-код для публикации маршрута & Запланировано \\
\hline
FR-34 & Встраиваемый iframe-виджет & Запланировано \\
\hline
FR-37 & Сезонность маршрутов и фильтрация & Запланировано \\
\hline
\end{tabularx}
\end{table}

\section{Анализ производительности}

Нефункциональные требования к производительности предусматривают время отклика API не более 500~мс на 95-м перцентиле при нагрузке 100 одновременных запросов, а также время обработки фотографии не более 10~с.

Для ключевых эндпоинтов сервиса маршрутов были измерены времена отклика при нагрузке. Запросы чтения (\texttt{GET /api/v1/routes/explore}, \texttt{GET /api/v1/shared/\{token\}}) показали медианное время отклика менее 30~мс благодаря PostgreSQL-индексам на полях \texttt{share\_token} и \texttt{user\_id}. Запросы записи (\texttt{POST /api/v1/routes}) занимают в среднем 45--80~мс, включая INSERT в PostgreSQL и публикацию задачи в NATS.

Время полного цикла обработки фотографии (от загрузки до WebSocket-уведомления) составляет 2--4~секунды для изображений размером до 5~МБ, что значительно ниже требуемого порога в 10~с. Основную часть этого времени занимает масштабирование через фильтр Ланцоша: алгоритм обеспечивает высокое качество за счёт вычислительных затрат. Для изображений с шириной менее 1920~пикселей масштабирование пропускается, что снижает время обработки до 0,5--1~с.

Все нефункциональные требования по производительности (NFR-01, NFR-02) выполнены. Доступность системы в Kubernetes-кластере обеспечивается за счёт механизма liveness/readiness проб и автоматического перезапуска упавших подов.

\section{Экономическое обоснование}

\subsection{Оценка трудоёмкости разработки}

Трудоёмкость разработки системы Guide Helper оценивалась экспертным методом по
этапам жизненного цикла разработки программного обеспечения. Результаты оценки
приведены в таблице~\ref{tab:labor}.

\begin{table}[H]
\caption{Трудоёмкость разработки системы по этапам}
\label{tab:labor}
\begin{tabularx}{\textwidth}{|l|X|c|}
\hline
\textbf{№} & \textbf{Этап разработки} & \textbf{Трудоёмкость, чел.-ч} \\
\hline
1 & Анализ предметной области и формулировка требований & 48 \\
\hline
2 & Проектирование архитектуры и схемы базы данных & 52 \\
\hline
3 & Реализация сервиса аутентификации & 40 \\
\hline
4 & Реализация сервиса маршрутов (основной API) & 80 \\
\hline
5 & Реализация конвейера обработки фотографий & 60 \\
\hline
6 & Интеграция ИИ-ассистента & 45 \\
\hline
7 & Разработка клиентского веб-приложения (React) & 100 \\
\hline
8 & Реализация вспомогательных сервисов (кеш, тайлы) & 30 \\
\hline
9 & Настройка CI/CD-пайплайна и Kubernetes & 35 \\
\hline
10 & Тестирование (модульное и функциональное) & 50 \\
\hline
11 & Разработка технической документации & 30 \\
\hline
\multicolumn{2}{|l|}{\textbf{Итого}} & \textbf{570} \\
\hline
\end{tabularx}
\end{table}

Суммарная трудоёмкость разработки составляет \textbf{570 человеко-часов}, что
соответствует примерно 3,5~месяцам работы одного разработчика при стандартной
занятости 160~часов в месяц.

\subsection{Расчёт стоимости разработки}

Расчёт стоимости разработки выполнен исходя из средней рыночной ставки разработчика
уровня junior/middle в России в 2025~году. Результаты расчёта приведены в
таблице~\ref{tab:cost}.

\begin{table}[H]
\caption{Расчёт стоимости разработки}
\label{tab:cost}
\begin{tabularx}{\textwidth}{|l|X|r|}
\hline
\textbf{№} & \textbf{Статья затрат} & \textbf{Сумма, руб.} \\
\hline
1 & Фонд оплаты труда (570~ч $\times$ 2\,200~руб./ч) & 1\,254\,000 \\
\hline
2 & Страховые взносы (30,2\% от ФОТ) & 378\,708 \\
\hline
3 & Накладные расходы (20\% от ФОТ) & 250\,800 \\
\hline
4 & Амортизация оборудования (ноутбук, 3~года, 4~мес.) & 14\,000 \\
\hline
\multicolumn{2}{|l|}{\textbf{Итого себестоимость разработки}} & \textbf{1\,897\,508} \\
\hline
\end{tabularx}
\end{table}

Таким образом, себестоимость разработки системы Guide Helper составляет
около \textbf{1\,898\,000~рублей}.

\subsection{Лицензирование и стоимость владения}

Все программные компоненты, использованные при разработке системы, распространяются
по свободным лицензиям, не требующим лицензионных отчислений: язык Rust
(MIT~/ Apache~2.0), язык Go (BSD-3-Clause), библиотека React (MIT), СУБД
PostgreSQL (PostgreSQL License), брокер NATS (Apache~2.0), объектное хранилище
MinIO (AGPL-3.0 для серверного развёртывания). Это снижает себестоимость
разработки и сопровождения платформы. Исходный код системы Guide Helper не
предполагается публиковать в открытом доступе: система рассматривается как
проприетарный программный продукт, предоставляемый конечным пользователям в виде
централизованного веб-сервиса. Такой подход сохраняет возможность коммерциализации,
управляемого развития и контроля над функциональностью платформы.

Годовые эксплуатационные расходы на инфраструктуру приведены в
таблице~\ref{tab:infra}.

\begin{table}[H]
\caption{Годовые эксплуатационные расходы на инфраструктуру}
\label{tab:infra}
\begin{tabularx}{\textwidth}{|l|X|r|}
\hline
\textbf{№} & \textbf{Статья расходов} & \textbf{Сумма, руб./год} \\
\hline
1 & Аренда VPS для Kubernetes-кластера (3~узла) & 72\,000 \\
\hline
2 & Доменное имя и DNS & 1\,500 \\
\hline
3 & Резервное копирование данных & 6\,000 \\
\hline
\multicolumn{2}{|l|}{\textbf{Итого}} & \textbf{79\,500} \\
\hline
\end{tabularx}
\end{table}

\subsection{Оценка экономического эффекта}

Экономический эффект от внедрения системы Guide Helper определяется снижением
временны́х затрат гидов и инструкторов на подготовку и передачу туристических
маршрутов клиентам. До внедрения системы типичный рабочий процесс включал: ручную
сортировку фотографий, формирование PDF-отчёта со скриншотами карты, рассылку
в мессенджерах. В среднем этот процесс занимал около 3~часов на один маршрут.

С использованием Guide Helper подготовка и публикация маршрута с геопривязанными
фотографиями занимает около 30~минут. Экономия времени на один маршрут составляет
2,5~часа. При среднем темпе работы гида 2~маршрута в месяц и ставке труда
2\,000~руб./ч годовая экономия для одного гида составит:

\[
  E = 2{,}5 \; \text{ч} \times 2 \; \text{маршрута/мес} \times 12 \; \text{мес}
      \times 2\,000 \; \text{руб./ч} = 120\,000 \; \text{руб./год}
\]

При использовании системы аудиторией из 20~гидов годовой экономический эффект
составит $20 \times 120\,000 = 2\,400\,000$~руб. Срок окупаемости разработки в
этом случае составит менее одного года:

\[
  T = \frac{1\,898\,000}{2\,400\,000} \approx 0{,}8 \; \text{года}
\]

Помимо прямой экономии времени, система обеспечивает дополнительные выгоды:
снижение зависимости от коммерческих платных платформ (экономия до 540\,000~руб./год
для организации с 10~гидами при сравнении с тарифами аналогичных зарубежных
сервисов), централизованный контроль над данными маршрутов и возможность адаптации
системы под корпоративные требования силами разработчика.

\section{Сценарий видеоролика для защиты}

Для защиты выпускной квалификационной работы подготовлен сценарий демонстрационного видеоролика продолжительностью 3~минуты. Сценарий построен по принципу решения проблемных кейсов и охватывает все три роли системы.

Сегмент 1 [0:00--0:25]. Постановка проблемы.\\
Представляется типичная проблема гида: как передать клиентам фотографии ключевых точек маршрута с географическим контекстом. Показаны альтернативы~--- PDF-файл, мессенджер~--- и их ограничения.

Сегмент 2 [0:25--1:25]. Роль «Гид/инструктор».\\
Демонстрируется: вход в приложение; создание маршрута с расстановкой 4~точек по карте Москвы; автоматическое построение маршрута через OSRM; загрузка фотографии с GPS-метками в EXIF; появление миниатюры на маркере после WebSocket-уведомления; панель статистики (расстояние, набор высоты, сложность); получение ссылки шеринга.

Сегмент 3 [1:25--2:10]. Роль «Турист».\\
Демонстрируется: открытие ссылки без авторизации; просмотр интерактивной карты с маркерами фото; просмотр прогноза погоды и профиля высот; регистрация; выставление оценки и написание комментария; использование каталога с поиском и фильтрацией.

Сегмент 4 [2:10--2:45]. Роль «Администратор».\\
Демонстрируется: административная панель; управление пользователями и ролями; модерация контента; создание категории; просмотр статистики платформы.

Сегмент 5 [2:45--3:00]. Архитектура системы.\\
Кратко демонстрируется диаграмма контейнеров системы; упоминаются ключевые технологии (Rust, Go, React, Kubernetes); приводится адрес демонстрационного стенда системы.

\section*{Выводы по третьей главе}
\addcontentsline{toc}{section}{Выводы по третьей главе}

По результатам тестирования системы Guide Helper подтверждено выполнение основных
функциональных и нефункциональных требований. Разработано 115~модульных тестов,
покрывающих доменную модель, JWT-сервис, бизнес-логику маршрутов, импорт GeoJSON,
HTTP-обработчики и ИИ-чат. Все тесты проходят в CI/CD-пайплайне при каждом коммите.

Функциональное тестирование по сценариям критического пути для трёх ролей
пользователей (гид, турист, администратор) подтвердило корректность реализации
33~из~37~требований (89\%). Четыре требования~--- AI-генерация описания,
QR-код, встраиваемый виджет и сезонность~--- находятся в разработке или запланированы.

Измеренное время отклика запросов чтения составляет менее 30~мс, запросов
записи~--- 45--80~мс; полный цикл обработки фотографии занимает 2--4~секунды.
Все нефункциональные требования по производительности (NFR-01, NFR-02) выполнены
с существенным запасом относительно установленных порогов в 500~мс и 10~секунд
соответственно.
